\section{Related Work}
\label{sec:related_work}

\subsection{Static Model Merging}
Model merging has emerged as a key pathway to combine the specialized capabilities of distinct expert neural networks without undergoing the high computational cost of joint training or multi-task fine-tuning~\cite{wortsman2022model, jin2022dataless}. Early and influential work in this area includes Task Arithmetic~\cite{ilharco2022editing}, which constructs task vectors by subtracting pre-trained base model weights from converging fine-tuned weights, allowing linear addition or subtraction of capabilities directly in the weight space. Model Soups~\cite{wortsman2022model} average weights of models fine-tuned with different hyperparameters to improve generalization. Fisher Weighted Averaging~\cite{matena2021merging} utilizes diagonal Fisher information matrices to weigh parameters according to their importance on the source tasks. To mitigate destructive parameter interference (where expert networks have conflicting sign values or magnitude imbalances at the same coordinates), TIES-merging~\cite{yadav2023ties} resolves conflicts by applying sign consensus and pruning redundant parameter updates. While these static merging strategies are training-free and fast to execute, they suffer from a fundamental limitation: they use fixed global coefficients across the entire network, failing to handle complex layer-wise and input-dependent variations in task conflicts.

\subsection{Active and Test-Time Model Merging}
To overcome the limitations of static averaging, active model merging paradigms optimize merging coefficients dynamically. A prominent example is AdaMerging~\cite{adamerging}, which introduces learnable merging coefficients. AdaMerging can be trained in either a supervised manner on training data or an unsupervised manner during test-time adaptation (TTA) on unlabeled test streams by minimizing prediction entropy. However, unconstrained coefficient adaptation on tiny online batches often leads to the Overfitting-Optimizer Paradox~\cite{regcalmerge, polymerge}. To regularize this adaptation loop, RegCalMerge~\cite{regcalmerge} introduces complex structural normalization terms and an Elastic Spatial Regularizer that imposes a quadratic distance penalty on parameter movements. PolyMerge~\cite{polymerge} restricts the learnable coefficients to low-degree polynomial trajectories over the depth of the model, reducing the parameter search space. Similarly, recent wave-inspired dynamic routing frameworks use complex cosine wave interference equations to smooth optimization landscapes. In contrast to these highly complex, multi-hyperparameter spatial regularizers and geometrical restrictions, PG-Merge demonstrates that a simple, training-free, non-parametric sparse gradient mask is fully sufficient to stabilize test-time coefficient optimization.

\subsection{Test-Time Adaptation (TTA)}
Test-time adaptation refers to the process of updating a pre-trained model on unlabeled test-time data to mitigate distribution shifts or customize the model on-the-fly~\cite{wang2020tent, liao2021are}. Tent~\cite{wang2020tent} popularized this paradigm by minimizing the Shannon entropy of model predictions across batch normalization parameters. While highly effective, TTA on low-sample test streams is notoriously unstable and susceptible to transductive collapse, where the network collapses to a trivial, constant prediction or overfits to local spurious correlations in the test batch~\cite{liang2023source, liao2021are}. Various works have sought to stabilize TTA using strong priors, contrastive losses, or temporal consistency. PG-Merge draws inspiration from these stability concerns, applying them to the specific domain of parameter-space model fusion. We demonstrate that restricting the active parameter updates to only a sparse subset of highly sensitive merging coefficients acts as a natural stabilizer for the TTA loop.

\subsection{Gradient Pruning and Sparse Optimization}
Our proposed technique is closely related to parameter-efficient fine-tuning (PEFT) and gradient selection methods in deep networks. PEFT methods, such as LoRA~\cite{hu2021lora}, Adapters~\cite{houlsby2019parameter}, and prefix tuning~\cite{li2021prefix}, restrict parameter updates to a small fraction of the model to prevent catastrophic forgetting and reduce memory requirements. In the context of optimization, gradient pruning and coordinate selection (e.g., MECTA) have been used to select only high-magnitude gradient components to update during fine-tuning. However, prior gradient selection methods have primarily focused on reducing backpropagation memory or speeding up training in standard supervised regimes. In this work, we repurpose gradient-magnitude sorting and sparse masking as an analytical, training-free regularizer inside the test-time model merging adaptation loop, showing that it naturally mitigates the Overfitting-Optimizer Paradox.